













\subsectionoptional{行列式的存在性}
\textit{本小节给出前一小节中
两个结论的证明。
本小节为选读内容。
这里展开的材料只会在
若尔当标准形小节中用到，而该小节也是选读的。}

我们想要证明：对任意阶数~$n$，
$\nbyn{n}$矩阵上的行列式函数是良定义的。
前一小节给出了
置换展开公式。\index{determinant!permutation expansion}%
\index{permutation expansion}
\begin{align*}
   %\deter{T}
   %=
   \begin{vmat}
      t_{1,1}  &t_{1,2}  &\ldots  &t_{1,n}  \\
      t_{2,1}  &t_{2,2}  &\ldots  &t_{2,n}  \\
               &\vdots                      \\
      t_{n,1}  &t_{n,2}  &\ldots  &t_{n,n}
   \end{vmat}
   &=
   \begin{array}[t]{l}
      t_{1,\phi_1(1)}t_{2,\phi_1(2)}\cdots
           t_{n,\phi_1(n)}\deter{P_{\phi_1}}       \\[.75ex]
      \>\hbox{}+t_{1,\phi_2(1)}t_{2,\phi_2(2)}\cdots
           t_{n,\phi_2(n)}\deter{P_{\phi_2}}       \\[.75ex]
      \>\alignedvdots                              \\
      \>\hbox{}+t_{1,\phi_k(1)}t_{2,\phi_k(2)}\cdots
           t_{n,\phi_k(n)}\deter{P_{\phi_k}} 
   \end{array}                                                 \\[1ex]
   &=\!\!\sum_{\text{permutations\ }\phi}\!\!\!\!
     t_{1,\phi(1)}t_{2,\phi(2)}\cdots t_{n,\phi(n)}
                                 \deter{P_{\phi}}
\end{align*}
这就把证明行列式是良定义的这一
问题化归为
只需证明行列式在置换矩阵的集合上
是良定义的。

置换矩阵可以通过行对换变成单位矩阵。
因此我们可以用记录对换次数的方法
来计算它的行列式。
然而，
我们仍必须证明该结果是良定义的。
回想一下困难所在：~矩阵
\begin{equation*}
   P_{\phi}=
   \begin{mat}[r]
      0  &1  &0  &0 \\
      1  &0  &0  &0 \\
      0  &0  &1  &0 \\
      0  &0  &0  &1
   \end{mat}
\end{equation*}
的行列式既可以用一次对换
\begin{equation*}
   P_{\phi}
   \grstep{\rho_1\leftrightarrow\rho_2}
   \begin{mat}[r]
      1  &0  &0  &0 \\
      0  &1  &0  &0 \\
      0  &0  &1  &0 \\
      0  &0  &0  &1
   \end{mat}
\end{equation*}
也可以用三次对换来计算。
\begin{equation*}
   P_{\phi}
   \grstep{\rho_3\leftrightarrow\rho_1}
   % \begin{mat}[r]
   %    0  &0  &1  &0 \\
   %    1  &0  &0  &0 \\
   %    0  &1  &0  &0 \\
   %    0  &0  &0  &1
   % \end{mat}
   \repeatedgrstep{\rho_2\leftrightarrow\rho_3}
   % \begin{mat}[r]
   %    0  &0  &1  &0 \\
   %    0  &1  &0  &0 \\
   %    1  &0  &0  &0 \\
   %    0  &0  &0  &1
   % \end{mat}
   \repeatedgrstep{\rho_1\leftrightarrow\rho_3}
   \begin{mat}[r]
      1  &0  &0  &0 \\
      0  &1  &0  &0 \\
      0  &0  &1  &0 \\
      0  &0  &0  &1
   \end{mat}
\end{equation*}
这两种化简所用的对换次数都是奇数，所以这里我们认为
\( \deter{P_{\phi}}=-1 \)；
但如果有某种方法用偶数次对换来实现，
那么同一个输入就会使行列式给出
两个不同的输出。
下面，\nearbycorollary{cor:ParityInversEqParitySwaps} 将证明
这种情况不会发生\Dash
不存在这样的置换矩阵，它能以两种方式经行对换变成单位矩阵，
其中一种用偶数次对换，另一种用奇数次对换。

% So the critical step will be a way to calculate whether the 
% number of swaps that it takes could be even or odd.

\begin{definition} \label{df:Inversion}
%<*df:Inversion>
在置换 $\phi=\sequence{\ldots,k,\ldots,j,\ldots}$ 中，
满足 $k>j$ 的元素相对于其自然顺序
构成一个\definend{逆序}。
类似地，在置换矩阵中，两行
\begin{equation*}
  P_{\phi}=
  \begin{mat}
    \vdots          \\
    \iota_{k} \\
    \vdots          \\
    \iota_{j} \\
    \vdots
  \end{mat}
\end{equation*}
若满足 \( k>j \)，则处于一个\definend{逆序}中。\index{inversion}%
\index{permutation!inversions}
%</df:Inversion>
\end{definition}

\begin{example}
这个置换矩阵
\begin{equation*}
  \begin{mat}[r]
    1  &0  &0  &0 \\
    0  &0  &1  &0 \\
    0  &1  &0  &0 \\
    0  &0  &0  &1
  \end{mat}
  =
  \begin{mat}
    \iota_1  \\
    \iota_3  \\
    \iota_2  \\
    \iota_4
  \end{mat}
\end{equation*}
只有一个逆序，即 \( \iota_3 \) 位于 \( \iota_2 \) 之前。
\end{example}

\begin{example} \label{ex:ThreeInversions}
这里有三个逆序：
\begin{equation*}
  \begin{mat}[r]
    0  &0  &1  \\
    0  &1  &0  \\
    1  &0  &0
  \end{mat}
  =
  \begin{mat}
    \iota_3  \\
    \iota_2  \\
    \iota_1
  \end{mat}
\end{equation*}
\( \iota_3 \) 位于 \( \iota_1 \) 之前，
\( \iota_3 \) 位于 \( \iota_2 \) 之前，且 \( \iota_2 \) 位于
\( \iota_1 \) 之前。
\end{example}

\begin{lemma} \label{le:SwapsChangeSgn}
%<*lm:SwapsChangeSgn>
置换矩阵中的一次行对换会使逆序的个数
从偶数变为奇数，或从奇数变为偶数。
%</lm:SwapsChangeSgn>
\end{lemma}

\begin{proof}
%<*pf:SwapsChangeSgn0>
考虑对第 $j$ 行与第 $k$ 行的对换，其中 $k>j$。

若这两行相邻
\begin{equation*}
   P_{\phi}=
   \begin{mat}
     \vdots           \\
     \iota_{\phi(j)}  \\
     \iota_{\phi(k)}  \\
     \vdots
   \end{mat}
 \grstep{\rho_k\leftrightarrow\rho_j}
   \begin{mat}
     \vdots           \\
     \iota_{\phi(k)}  \\
     \iota_{\phi(j)}  \\
     \vdots
   \end{mat}
\end{equation*}
那么由于涉及这对之外其他行的逆序不受影响，
这次对换使逆序的总数改变一，
究竟是消去一个逆序还是产生一个逆序，
取决于 \( \phi(j)>\phi(k) \) 是否成立。
因此，逆序的总数从奇数变为偶数，
或从偶数变为奇数。
%</pf:SwapsChangeSgn0>

%<*pf:SwapsChangeSgn1>
若这两行不相邻，则我们可以通过一系列相邻行的对换来交换它们，
先把第~$k$ 行向上移
\begin{equation*}
   \begin{mat}
     \vdots             \\
     \iota_{\phi(j)}    \\
     \iota_{\phi(j+1)}  \\
     \iota_{\phi(j+2)}  \\
     \vdots             \\
     \iota_{\phi(k)}    \\
     \vdots
   \end{mat}
  \grstep{\rho_k\swap\rho_{k-1}}
  \repeatedgrstep{\rho_{k-1}\swap\rho_{k-2}}
  \cdots
  \grstep{\rho_{j+1}\swap\rho_j}
   \begin{mat}
     \vdots             \\
     \iota_{\phi(k)}    \\
     \iota_{\phi(j)}    \\
     \iota_{\phi(j+1)}  \\
     \vdots             \\
     \iota_{\phi(k-1)}  \\
     \vdots
   \end{mat}
\end{equation*}
%</pf:SwapsChangeSgn1>
%<*pf:SwapsChangeSgn2>
再把第~$j$ 行向下移。
\begin{equation*}
  \grstep{\rho_{j+1}\swap\rho_{j+2}}
  \repeatedgrstep{\rho_{j+2}\swap\rho_{j+3}}
  \cdots
  \grstep{\rho_{k-1}\swap\rho_k}
   \begin{mat}
     \vdots             \\
     \iota_{\phi(k)}    \\
     \iota_{\phi(j+1)}  \\
     \iota_{\phi(j+2)}  \\
     \vdots             \\
     \iota_{\phi(j)}    \\
     \vdots
   \end{mat}
\end{equation*}
这些相邻对换中的每一次都使逆序的个数
从奇数变为偶数或从偶数变为奇数。
对换的总次数 \( (k-j)+(k-j-1) \) 是奇数。
于是，总体上逆序的个数
从偶数变为奇数，或从奇数变为偶数。
%</pf:SwapsChangeSgn2>
\end{proof}

\begin{corollary}  \label{cor:ParityInversEqParitySwaps}
%<*co:ParityInversEqParitySwaps>
若置换矩阵有奇数个逆序，则把它换成
单位矩阵需要奇数次对换。
若它有偶数个逆序，则换成单位矩阵需要偶数次对换。
%</co:ParityInversEqParitySwaps>
\end{corollary}

\begin{proof}
%<*pf:ParityInversEqParitySwaps>
单位矩阵有零个逆序。
把奇数变成零需要奇数次对换，
把偶数变成零需要偶数次对换。
%</pf:ParityInversEqParitySwaps>
\end{proof}

\begin{example}
\nearbyexample{ex:ThreeInversions} 中的矩阵
用一次对换 $\rho_1\leftrightarrow\rho_3$ 就能变成单位矩阵。
（所以对换的次数不必与逆序的个数相同，
但这两个数的奇偶性是一样的。）
\end{example}

\begin{definition} \label{df:Signum}
%<*df:Signum>
置换的\definend{符号}\index{permutation!signum}\index{signum}%
\index{sgn!see {signum}}
\( \sgn(\phi) \)
在 \( \phi \) 中逆序个数为奇数时是 \( -1 \)，为偶数时是
\( +1 \)。
%</df:Signum>
\end{definition}

\begin{example}
使用 \nearbyexample{ex:AllTwoThreePerms} 中
$3$ 元置换的记号，我们有
\begin{equation*}
  P_{\phi_1}=
  \begin{mat}
    1  &0  &0  \\
    0  &1  &0  \\
    0  &0  &1
  \end{mat}
  \qquad
  P_{\phi_2}=
  \begin{mat}
    1  &0  &0  \\
    0  &0  &1  \\
    0  &1  &0  
  \end{mat}
\end{equation*}
所以
\( \sgn(\phi_1)=1 \)，因为没有逆序，
而 \( \sgn(\phi_2)=-1 \)，因为有一个逆序。
\end{example}

我们仍未证明行列式函数是良定义的，
因为除了行对换之外，
我们还没有考虑置换矩阵上的其他行变换。
我们将巧妙地绕过这个问题。
%<*DefiningDFunction>
定义函数
\( \map{d}{\matspace_{\nbyn{n}}}{\Re} \)：
修改置换展开公式，
把 $\deter{P_\phi}$ 换成 $\sgn(\phi)$。
\begin{equation*}
  d(T)=
  \sum_{\text{permutations\ }\phi}t_{1,\phi(1)}t_{2,\phi(2)}\cdots t_{n,\phi(n)}
                                 \cdot\sgn(\phi)
\end{equation*}
% This gives the same value as the permutation expansion
% because the corollary shows that $\det(P_\phi)=\sgn(\phi)$.
这个公式的优点是逆序的个数
显然是良定义的\Dash  直接数出来即可。
因此，当我们证明了~$d$
满足行列式定义中的条件时，
也就证明了一个~$\nbyn{n}$阶行列式函数的
存在性。
%</DefiningDFunction>

\begin{lemma} \label{lm:DeterminantsExist}
\index{determinant!exists}
%<*lm:DeterminantsExist>
上面的函数 \( d \) 是一个行列式。
从而 $\det_{\nbyn{n}}$ 对每个 $n$ 都存在。
%</lm:DeterminantsExist>
\end{lemma}

\begin{proof}
%<*pf:DeterminantsExist0>
我们必须验证它满足行列式定义中的四个条件，
即 \nearbydefinition{def:Det}。
%</pf:DeterminantsExist0>

%<*pf:DeterminantsExist1>
条件~(4) 容易验证：设 $I$ 是 $\nbyn{n}$ 阶单位矩阵，
在
\begin{equation*}
  d(I)=
  \sum_{\text{perm\ }\phi}
    \iota_{1,\phi(1)}\iota_{2,\phi(2)}\cdots \iota_{n,\phi(n)}
                                 \sgn(\phi)
\end{equation*}
中，求和式里所有的项都为零，
只有置换~$\phi$ 为恒等的那一项例外，
这一项给出沿对角线的乘积，
而该乘积为一。
%</pf:DeterminantsExist1>

%<*pf:DeterminantsExist2>
对于条件~(3)，设 \( T\smash[b]{\grstep{k\rho_i}}\hat{T} \)，
考虑 $d(\hat{T})$。
\begin{multline*}
  \sum_{\text{perm\ }\phi}\!\!
    \hat{t}_{1,\phi(1)}
     \cdots\hat{t}_{i,\phi(i)}\cdots\hat{t}_{n,\phi(n)}
                                 \sgn(\phi)   \\
  =\sum_{\phi}
     t_{1,\phi(1)}\cdots kt_{i,\phi(i)}\cdots t_{n,\phi(n)}
                                 \sgn(\phi)
\end{multline*}
提出因子~\( k \) 便得到所要的等式。
\begin{equation*}
  =k\cdot\sum_{\phi}
     t_{1,\phi(1)}\cdots t_{i,\phi(i)}\cdots t_{n,\phi(n)}
                                 \sgn(\phi)                 
  =k\cdot d(T)
\end{equation*}
%</pf:DeterminantsExist2>

%<*pf:DeterminantsExist3>
对于~(2)，设
\( T\smash[b]{\grstep{\rho_i\swap\rho_j}}\hat{T} \)。
我们必须证明 $d(\hat{T})$ 是 $d(T)$ 的相反数。
\begin{equation*}
  d(\hat{T})=
  \sum_{\text{perm\ }\phi}\!\!
    \hat{t}_{1,\phi(1)}
    \cdots\hat{t}_{i,\phi(i)}
    \cdots\hat{t}_{j,\phi(j)}
    \cdots \hat{t}_{n,\phi(n)}
                                 \sgn(\phi)
    \tag{$*$}
\end{equation*}
我们将证明（$*$）中的每一项都对应 $d(T)$ 中的一项，
且这两项互为相反数。
考虑 $d(\hat{T})$ 的多重线性展开式中给出项
    $\hat{t}_{1,\phi(1)}
    \cdots\hat{t}_{i,\phi(i)}
    \cdots\hat{t}_{j,\phi(j)}
    \cdots \hat{t}_{n,\phi(n)}
                                  \sgn(\phi)$ 的矩阵。
\begin{equation*}
  \begin{mat}
     \ &                  &\vdots &               &     \\
       &\hat{t}_{i,\phi(i)}  &       &               &    \\
       &                &\vdots &                 &  \\
       &                &       &\hat{t}_{j,\phi(j)} &  \\
       &                &\vdots                   &
  \end{mat}
\end{equation*}
%</pf:DeterminantsExist3>
%<*pf:DeterminantsExist4>
它是对矩阵
施加 $\rho_i\swap\rho_j$ 操作的结果。
\begin{equation*}
  \begin{mat}
       \ &              &\vdots &           &        \\
         &              &       &t_{i,\phi(j)} &  \\
         &              &\vdots &           &       \\
         &t_{j,\phi(i)}    &       &          &        \\
         &              &\vdots &          &
  \end{mat}
\end{equation*}
也就是说，带帽 $t$ 的项对应于来自 $d(T)$ 展开式的这一项：
 \(    t_{1,\sigma(1)}
       \cdots t_{j,\sigma(j)}
       \cdots t_{i,\sigma(i)}
       \cdots t_{n,\sigma(n)} \sgn(\sigma)
       \)，其中置换
\( \sigma \)
与 \( \phi \) 相同，只是把第 \( i \) 个数与第 \( j \) 个数互换：
\( \sigma(i)=\phi(j) \) 且 \( \sigma(j)=\phi(i) \)。
这两项的各因子相同：
$\hat{t}_{1,\phi(1)}=t_{1,\sigma(1)}$，\ldots，
包括来自对换过的两行的元素
$\hat{t}_{i,\phi(i)}=t_{j,\phi(i)}=t_{j,\sigma(j)}$
和 $\hat{t}_{j,\phi(j)}=t_{i,\phi(j)}=t_{i,\sigma(i)}$。 
但这两项互为相反数，
因为由 \nearbylemma{le:SwapsChangeSgn} 有 \( \sgn(\phi)=-\sgn(\sigma) \)。
%</pf:DeterminantsExist4>

%<*pf:DeterminantsExist5>
现在，任意置换 $\phi$ 都可以由某个
其他置换 $\sigma$ 通过这样的对换得到，
而且方式有且仅有一种。
因此（$*$）中的求和实际上是对所有置换的求和，
每个置换恰好取一次。
\begin{align*}
  d(\hat{T})
  &=
  \sum_{\text{perm\ }\phi}\!\!
    \hat{t}_{1,\phi(1)}
    \cdots\hat{t}_{i,\phi(i)}
    \cdots\hat{t}_{j,\phi(j)}
    \cdots \hat{t}_{n,\phi(n)}
                                 \sgn(\phi)   \\
  &=\sum_{\text{perm\ }\sigma}\!\!
     t_{1,\sigma(1)}
    \cdots t_{j,\sigma(j)}
    \cdots t_{i,\sigma(i)}
    \cdots t_{n,\sigma(n)}
                                 \cdot\bigl(-\sgn(\sigma)\bigr)        
\end{align*}
于是 \( d(\hat{T})=-d(T) \)。
%</pf:DeterminantsExist5>

%<*pf:DeterminantsExist6>
最后，对于条件~(1)，设
\( T\smash[b]{\grstep{k\rho_i+\rho_j}}\hat{T} \)。
\begin{align*}
  d(\hat{T})
  &=\sum_{\text{perm\ }\phi}
    \hat{t}_{1,\phi(1)}\cdots\hat{t}_{i,\phi(i)}
    \cdots\hat{t}_{j,\phi(j)}\cdots\hat{t}_{n,\phi(n)}
                                 \sgn(\phi)                \\
  &=\sum_{\phi}
    t_{1,\phi(1)}\cdots t_{i,\phi(i)}
    \cdots (kt_{i,\phi(j)}+t_{j,\phi(j)})\cdots t_{n,\phi(n)}
                                 \sgn(\phi)
\end{align*}
% (notice:~that's 
% \( kt_{i,\phi(j)} \), not \( kt_{j,\phi(j)} \)).
对 $kt_{i,\phi(j)}+t_{j,\phi(j)}$ 中的加法使用分配律。
\begin{equation*}
  = 
      \begin{aligned}[t]
        &\smash[b]{\sum_{\smash{\phi}}}\big[t_{1,\phi(1)}\cdots t_{i,\phi(i)}
          \cdots kt_{i,\phi(j)}\cdots t_{n,\phi(n)}
                                 \sgn(\phi)         \\ % [-1ex]
        &\hbox{}\qquad\quad\hbox{}
           +t_{1,\phi(1)}\cdots t_{i,\phi(i)}
           \cdots t_{j,\phi(j)}\cdots t_{n,\phi(n)}
                                 \sgn(\phi)\big]
      \end{aligned}                                               
\end{equation*}
把它拆成两个和式。
\begin{equation*}
  =\begin{aligned}[t] 
         &\smash[b]{\sum_{\smash{\phi}}} \:t_{1,\phi(1)}\cdots t_{i,\phi(i)}
            \cdots kt_{i,\phi(j)}\cdots t_{n,\phi(n)}
                                 \sgn(\phi)           \\ % [-1ex]               
         &\hbox{}\qquad\quad\hbox{}
            +\smash[b]{\sum_{\phi}}\:
            t_{1,\phi(1)}\cdots t_{i,\phi(i)}
            \cdots t_{j,\phi(j)}\cdots t_{n,\phi(n)}
                                 \sgn(\phi)        
     \end{aligned}                                               
\end{equation*}
识别出第二个和式。
\begin{equation*}
  =\begin{aligned}[t]
       &k\cdot \smash[b]{\sum_{\smash{\phi}}}\:
            t_{1,\phi(1)}\cdots t_{i,\phi(i)}
            \cdots t_{i,\phi(j)}\cdots t_{n,\phi(n)}
                                 \sgn(\phi)          \\  % [-.75ex]
       &\hbox{}\qquad\quad\hbox{}
             +d(T)          
     \end{aligned}
\end{equation*}
%</pf:DeterminantsExist6>
%<*pf:DeterminantsExist7>
考虑各项
\( t_{1,\phi(1)}\cdots t_{i,\phi(i)} \cdots t_{i,\phi(j)}\cdots t_{n,\phi(n)}
   \sgn(\phi)  \)。
注意下标；元素是 \( t_{i,\phi(j)} \)，而不是 \( t_{j,\phi(j)} \)。
这些项的和是矩阵~\( S \) 的行列式，
它等于 \( T \)，
只是 $S$ 的第~$j$ 行是 $T$ 的第~$i$ 行的副本，
也就是说，$S$ 有两行相同。
用与证明 \nearbylemma{le:IdenRowsDetZero} 相同的方式可以看出
$d(S)=0$：交换 $S$ 的相同的两行会改变
$d(S)$ 的符号，但由于矩阵经这次交换并不改变，
$d(S)$ 的值也必须不变，因此这个值必为零。
%</pf:DeterminantsExist7>
\end{proof}

至此我们证明了每个阶数~$\nbyn{n}$ 的行列式函数都存在。
我们已经知道，对每个阶数，行列式至多有一个。
因此，对每个阶数，行列式函数有且仅有一个。

我们以证明来自前一小节的另一个尚待证明的结论
来结束本小节。

\begin{theorem} \label{th:DetMatrixEqualsDetTrans}
\index{transpose!determinant}
%<*th:DetMatrixEqualsDetTrans>
矩阵的行列式等于其转置的行列式。
%</th:DetMatrixEqualsDetTrans>
\end{theorem}

\begin{proof}
%<*pf:DetMatrixEqualsDetTrans0>
最好通过作出一般的
$\nbyn{3}$ 情形来理解这个证明。
该论证适用于 $\nbyn{n}$ 情形，这一点将是显然的。

比较矩阵~$T$ 的置换展开
\begin{align*}
  \begin{vmat}
    t_{1,1}  &t_{1,2}  &t_{1,3}  \\
    t_{2,1}  &t_{2,2}  &t_{2,3}  \\
    t_{3,1}  &t_{3,2}  &t_{3,3}  
  \end{vmat}
  &=
  t_{1,1}t_{2,2}t_{3,3}
  \begin{vmat}
    1 &0 &0 \\
    0 &1 &0 \\
    0 &0 &1
  \end{vmat}
  +t_{1,1}t_{2,3}t_{3,2}
  \begin{vmat}
    1 &0 &0 \\
    0 &0 &1 \\
    0 &1 &0
  \end{vmat}            \\
  &\quad+
  t_{1,2}t_{2,1}t_{3,3}
  \begin{vmat}
    0 &1 &0 \\
    1 &0 &0 \\
    0 &0 &1
  \end{vmat}
  +t_{1,2}t_{2,3}t_{3,1}
  \begin{vmat}
    0 &1 &0 \\
    0 &0 &1 \\
    1 &0 &0
  \end{vmat}           \\
  &\quad+
  t_{1,3}t_{2,1}t_{3,2}
  \begin{vmat}
    0 &0 &1 \\
    1 &0 &0 \\
    0 &1 &0
  \end{vmat}
  +t_{1,3}t_{2,2}t_{3,1}
  \begin{vmat}
    0 &0 &1 \\
    0 &1 &0 \\
    1 &0 &0
  \end{vmat}
\end{align*}
与其转置的置换展开。
%</pf:DetMatrixEqualsDetTrans0>
%<*pf:DetMatrixEqualsDetTrans1>
\begin{align*}
  \begin{vmat}
    t_{1,1}  &t_{2,1}  &t_{3,1}  \\
    t_{1,2}  &t_{2,2}  &t_{3,2}  \\
    t_{1,3}  &t_{2,3}  &t_{3,3}  
  \end{vmat}
  &=
  t_{1,1}t_{2,2}t_{3,3}
  \begin{vmat}
    1 &0 &0 \\
    0 &1 &0 \\
    0 &0 &1
  \end{vmat}
  +t_{1,1}t_{3,2}t_{2,3}
  \begin{vmat}
    1 &0 &0 \\
    0 &0 &1 \\
    0 &1 &0
  \end{vmat}            \\
  &\quad+
  t_{2,1}t_{1,2}t_{3,3}
  \begin{vmat}
    0 &1 &0 \\
    1 &0 &0 \\
    0 &0 &1
  \end{vmat}
  +t_{2,1}t_{3,2}t_{1,3}
  \begin{vmat}
    0 &1 &0 \\
    0 &0 &1 \\
    1 &0 &0
  \end{vmat}           \\
  &\quad+
  t_{3,1}t_{1,2}t_{2,3}
  \begin{vmat}
    0 &0 &1 \\
    1 &0 &0 \\
    0 &1 &0
  \end{vmat}
  +t_{3,1}t_{2,2}t_{1,3}
  \begin{vmat}
    0 &0 &1 \\
    0 &1 &0 \\
    1 &0 &0
  \end{vmat}
\end{align*}
先比较六个 $t$ 的乘积。
$T$ 展开式中的那些乘积与转置展开式中的相同；
例如，$t_{1,2}t_{2,3}t_{3,1}$ 在上面，
而 $t_{3,1}t_{1,2}t_{2,3}$ 在下面。
这完全合理\Dash 上面的六个来自从~$T$ 的每一行、每一列
各取一个元素的所有方式，
而下面的六个是从~$T$ 的每一列、每一行
各取一个元素的所有方式，
所以它们当然是同一个集合。
%</pf:DetMatrixEqualsDetTrans1>

%<*pf:DetMatrixEqualsDetTrans2>
接下来观察到，在这两个展开式中，
每个 $t$ 乘积表达式所对应的
置换矩阵不一定相同。
例如，在上面，$t_{1,2}t_{2,3}t_{3,1}$ 对应于映射
$1\mapsto 2$, $2\mapsto 3$, $3\mapsto 1$ 的矩阵。
在下面，$t_{3,1}t_{1,2}t_{2,3}$
对应于映射
$1\mapsto 3$, $2\mapsto 1$, $3\mapsto 2$ 的矩阵。
第二个映射是第一个的逆映射。
这也完全合理\Dash 两者
矩阵转置与逆映射都把 $1,2$ 换成 $2,1$，
把 $2,3$ 换成 $3,2$，并把 $3,1$ 换成 $1,3$。
 
最后我们指出，$P_{\phi}$ 的行列式等于
$P_{\phi^{-1}}$ 的行列式，
正如 \nearbyexercise{exer:PermInvFacts} %
所表明的。
%</pf:DetMatrixEqualsDetTrans2>
\end{proof}


